\hspace{3mm}

\centerline{\bf \Huge 9 класс}

\begin{flushright}
        \scriptsize{}
\end{flushright}

\problem{\textbf{1}}
В ряд стоят четыре человека, среди которых есть Маша (которая всегда говорит правду), Алина (которая всегда лжёт) и преподаватели Расул Владимирович с Анджуром Аркадьевичем (которые могут сказать как правду, так и солгать).

--- Самый левый сказал: <<Рядом со мной стоит Маша>>.

--- Второй слева сказал: <<Рядом со мной стоит Алина>>.

--- Третий слева сказал: <<Рядом со мной стоят преподаватели>>.

--- Самый правый сказал: <<Рядом со мной стоит Маша>>.

Определите, кто и где может стоять.

\problem{\textbf{2}}
Положительные числа $x,y,z$ таковы, что $xyz=1$. Чему может быть равно значение выражения:

$$
\left(x+\frac{1}{x}\right)^2+\left(y+\frac{1}{y}\right)^2+\left(z+\frac{1}{z}\right)^2-\left(x+\frac{1}{x}\right)\left(y+\frac{1}{y}\right)\left(z+\frac{1}{z}\right) ?
$$

%https://www.komal.hu/feladat?a=feladat&f=B5335&l=en

\problem{\textbf{3}}
Обозначим за $p(n)$ количество различных простых делителей $n>1$. Докажите, что существует бесконечно много натуральных $n$, для которых число $p(n+1)-p(n)$ чётно.


\problem{\textbf{4}}
В равнобедренном треугольнике $A B C\ (A B=A C)$ точка $M$ --- середина $B C$. На меньшей дуге $A M$ окружности $(A B M)$ выбрана произвольная точка $X$, а внутри угла $\angle B M A$ выбрана такая точка $T$, что $\angle T M X=90^{\circ}$ и $T X=B X$. Докажите, что разность углов $\angle M T B-\angle M T C$ не зависит от выбора точки $X$.

\problem{\textbf{5}}
На плоскости отмечено 36 точек, никакие три из которых не лежат на одной прямой. Некоторые пары отмеченных точек соединены отрезком так, что из каждой отмеченной точки выходит не более трёх отрезков. Какое наибольшее количество различных замкнутых 4-звенных ломаных может получиться? (Неважно, где у ломаной начало и как она ориентирована).